\begin{tabular}{lcc}
\hline
 & \textbf{Non-marginalized areas} & \textbf{Marginalized areas} \\
\hline

\multicolumn{3}{l}{\textbf{A. Progresa program intensity}} \\
\quad 1999 & 0.10 (0.13) & 0.44 (0.22) \\
\quad 2005 & 0.28 (0.19) & 0.72 (0.17) \\[6pt]

\multicolumn{3}{l}{\textbf{B. Characteristics of municipalities, 1990}} \\
\multicolumn{3}{l}{\quad \textit{Components of the marginality index}} \\
\quad Illiterate population & 13.33 (6.57) & 33.33 (13.3) \\
\quad With less than primary education & 45.74 (12.32) & 69.56 (9.68) \\
\quad Without a toilet & 25.94 (16.62) & 60.30 (18.54) \\
\quad Without electricity & 11.81 (10.29) & 36.39 (24.72) \\
\quad Without running water & 19.43 (14.87) & 50.65 (24.07) \\
\quad With crowding & 59.84 (9.79) & 73.96 (8.24) \\
\quad With dirt floor & 22.05 (14.21) & 61.98 (21.58) \\
\quad In communities with $<$5{,}000 inhabitants & 60.18 (36.15) & 95.50 (13.54) \\
\quad Earning $<$2× minimum wage & 69.19 (11.60) & 85.82 (8.55) \\[6pt]

Number of municipalities & 1,234 & 1,119 \\
\hline
\end{tabular}
